% Solution: Tadelis (2013) Exercise 11.5 "Asymmetric Patience 1". Input by \solitem.
% Reused from lista-reais-lectures-01-04-solutions (item 23), deepened. Check: check_C_L05.py.
\begin{exsol}{Tadelis Exercise 11.5}{236--237}{Exam level}{4}
\asked In three-period alternating-offer bargaining over a pie of 1, with player~1 proposing first
and discount factors $\delta_1 \ne \delta_2$: (a) the SPE outcome; (b) player~1's advantage when
$\delta_1 = \delta_2$; (c) when player~2 has the advantage, and why.

\method Backward induction in finite alternating-offer bargaining (Tadelis \S11.2, pp.~224--227;
Gibbons \S2.1.D, pp.~68--70). In the last period the proposer takes everything. In each earlier
period the proposer offers the responder exactly the discounted value of what the responder would
get next period, and the responder accepts at indifference (the argument of Prop.~11.2, p.~223;
agreement in round one, Prop.~11.3, p.~226). Slides: Lecture~4, slides~9--10 (backward induction)
and slide~17 (SPE). Lecture~5's bargaining part has no slides yet. Notation as in
\texttt{Notas/L04/L04-06-finite-alternating-offers.md}.

\solution
\textbf{Notation.} $x_t$ = player~1's share of the pie when agreement is reached in period
$t \in \{1,2,3\}$, so player~2 gets $1 - x_t$. Player~$i$'s payoff from agreement $(x, 1-x)$ in
period $t$ is $\delta_i^{t-1}$ times his share. Players~1, 2, 1 propose in periods 1, 2, 3, and no
agreement in period~3 gives $(0,0)$. ``Advantage'' means a share above $\frac12$.

\textbf{Step 1: period 3 (player~1 proposes, last period).} Player~2 compares accepting $1-x$ with
rejecting, which gives $0$. She accepts every $x < 1$ strictly. At $x = 1$ she is indifferent, and
in an SPE she must accept: if she rejected $x = 1$, player~1 would want the largest $x < 1$, which
does not exist, so he would have no best offer. Player~1 offers $x_3 = 1$. Period-3 values:
player~1 gets $1$, player~2 gets $0$; in period-1 terms, $(\delta_1^2,\ 0)$.

\textbf{Step 2: period 2 (player~2 proposes).} Player~1 accepts $x$ iff $x$ is at least what
rejecting gives him, namely $1$ one period later, worth $\delta_1 \cdot 1$ today:
\begin{equation}
  x \;\ge\; \delta_1\, x_3 = \delta_1 . \label{eq:tadelis-11-5-x2}
\end{equation}
Player~2 keeps the most she can by offering exactly $x_2 = \delta_1$ (by the same argument as in
Step~1, player~1 accepts at equality). She then gets $1 - \delta_1$. Making an offer player~1 rejects
gives her $\delta_2 \cdot 0 = 0 < 1 - \delta_1$, so she makes the acceptable offer.

\textbf{Step 3: period 1 (player~1 proposes).} Player~2 accepts $x$ iff her share now is at least
the value of rejecting, i.e.\ $1 - x_2$ one period later, discounted by her own $\delta_2$:
\[
  1 - x \;\ge\; \delta_2\,(1 - x_2) = \delta_2\,(1 - \delta_1) .
\]
Player~1 keeps as much as possible:
\begin{equation}
  x_1 = 1 - \delta_2(1 - \delta_1) = 1 - \delta_2 + \delta_1\delta_2 . \label{eq:tadelis-11-5-x1}
\end{equation}
He prefers this to a rejected offer, which gives him $\delta_1 x_2 = \delta_1^2$: indeed
$x_1 - \delta_1^2 = (1 - \delta_1)(1 + \delta_1 - \delta_2) > 0$.

\textbf{(a) The SPE.} Strategies at every node:
\begin{itemize}
  \item player~1: offer $x_1 = 1 - \delta_2(1-\delta_1)$ in period~1; in period~2 accept iff
  offered $x \ge \delta_1$; in period~3 offer $x_3 = 1$;
  \item player~2: in period~1 accept iff $1 - x \ge \delta_2(1-\delta_1)$; in period~2 offer
  $x_2 = \delta_1$; in period~3 accept every offer.
\end{itemize}
Each acceptance rule holds after every possible offer, including offers never made on the path, so
every subgame is played optimally. The outcome is immediate agreement:
\[
  (v_1, v_2) = \bigl(1 - \delta_2(1 - \delta_1),\ \ \delta_2(1 - \delta_1)\bigr).
\]
The tie-breaking argument of Step~1 applies at every round, so this SPE is unique.

\textbf{(b) Equal discount factors.} Setting $\delta_1 = \delta_2 = \delta$ in
\eqref{eq:tadelis-11-5-x1}:
\[
  v_1 = 1 - \delta + \delta^2, \qquad v_2 = \delta - \delta^2,
\]
\[
  v_1 - v_2 = 1 - 2\delta + 2\delta^2 = (1 - \delta)^2 + \delta^2 \;>\; 0 .
\]
Player~1 gets strictly more than half for every $\delta$. He proposes both first and last. The
edge $(1-\delta)^2 + \delta^2$ is smallest at $\delta = \frac12$ (derivative $-2(1-\delta) + 2\delta =
0$), where it equals $\frac12$ and $v_1 = \frac34$.

\textbf{(c) When player 2 is ahead.} Player~2 gets more than half iff
\[
  \delta_2\,(1 - \delta_1) > \tfrac12 \iff \delta_1 < 1 - \frac{1}{2\delta_2} .
\]
This needs $\delta_2 > \frac12$ (otherwise $\delta_2(1-\delta_1) \le \frac12$ for every $\delta_1$) and
$\delta_1$ low enough. Example: $\delta_2 = 0.9$ gives the cut-off $\delta_1 < 1 - \frac{1}{1.8} =
\frac49$; with $\delta_1 = 0.2$, $v_2 = 0.9 \times 0.8 = 0.72$.
\emph{Why.} Player~1's power comes from the last period. An impatient player~1 values it at only
$\delta_1$ from period~2's viewpoint, so player~2 keeps $1 - \delta_1$ when she proposes. A patient
player~2 values that period-2 position almost fully ($\delta_2$), so in period~1 player~1 must
concede $\delta_2(1 - \delta_1)$. Patience is bargaining power:
\[
  \frac{\partial v_1}{\partial \delta_1} = \delta_2 > 0, \qquad
  \frac{\partial v_1}{\partial \delta_2} = -(1 - \delta_1) < 0 .
\]
The infinite-horizon formula $v_1 = \frac{1 - \delta_2}{1 - \delta_1\delta_2}$ (Tadelis \S11.3;
\texttt{Notas/L04/L04-07}) has the same signs.

\interp First-mover and last-mover advantages both favour player~1, so only a large patience gap
in player~2's favour can overturn them. As $\delta \to 0$ (equal patience) the first-mover
advantage takes over and $v_1 \to 1$. As $\delta \to 1$ the last-mover advantage takes over and again
$v_1 \to 1$, because waiting for period~3 costs player~1 nothing. Patience matters because each
player's threat to reject is worth what the next round is worth \emph{to him}: the more patient
player loses less by waiting, so the other side must pay him more to settle now.

\verify \texttt{check\_C\_L05.py} redoes the recursion symbolically ($x_1 = 1 - \delta_2 +
\delta_1\delta_2$; $v_1 - v_2 = (1-\delta)^2 + \delta^2$ with minimum $\frac12$ at $\delta = \frac12$; limits
1 at both ends; the two derivatives). It tests the condition in (c) at the boundary $\delta_1 =
\frac49$ for $\delta_2 = \frac9{10}$ and at $\frac49 \pm \frac1{1000}$, confirms that player~2 can never be
ahead when $\delta_2 \le \frac12$ (grid of step $\frac1{100}$), and repeats the backward induction
numerically with offers on a grid of step $\frac1{1000}$ for three pairs; e.g.\ $(\delta_1,\delta_2) =
(0.9, 0.7)$ gives $x_1 = 0.93$. Output: ``Tadelis 11.5: v1 = 1 - d2(1 - d1), v2 = d2(1 - d1); equal
delta edge (1-d)\^{}2 + d\^{}2 >= 1/2; P2 ahead iff d2(1-d1) > 1/2 (boundary 4/9 at d2 = 9/10 checked
+-1/1000)''. Tadelis prints no answer.

\mistakes
\begin{itemize}
  \item Discounting player~2's continuation with $\delta_1$ (or with a common $\delta$). Each player
  values the future with \emph{his own} factor: player~2's period-1 threat is $\delta_2(1-\delta_1)$.
  \item Letting player~2 reject at indifference in period~3 (or player~1 in period~2). Then the
  proposer has no best offer and there is no SPE; acceptance at indifference is forced.
  \item Answering (b) with ``player~1 moves first'' only. With three periods he also moves last, which
  is why his edge does not vanish as $\delta \to 1$.
\end{itemize}

\answer{(a) Immediate agreement: player~1 gets $1 - \delta_2(1 - \delta_1)$, player~2 gets
$\delta_2(1 - \delta_1)$ (offers: $x_3 = 1$, $x_2 = \delta_1$, $x_1 = 1-\delta_2+\delta_1\delta_2$). (b) With
$\delta_1 = \delta_2 = \delta$: $v_1 - v_2 = (1-\delta)^2 + \delta^2 > 0$ (at least $\frac12$). (c) Player~2 is
ahead iff $\delta_2(1 - \delta_1) > \frac12$, which needs player~2 patient ($\delta_2 > \frac12$) and
player~1 impatient ($\delta_1 < 1 - \frac1{2\delta_2}$). Player~1's last-period power is worth little to an
impatient player~1, while player~2's period-2 power is worth a lot to a patient player~2.}
\end{exsol}
